%=====================================================================
%  RT-Comops Backend — Rapport d'Analyse & Conception
%  Module : auth-core (YowAuth)
%  Style aligné sur le Cahier de Conception RT-Comops (ConceptionKernel V3)
%  Compilable sur Overleaf (pdfLaTeX).  Encodage UTF-8.
%=====================================================================
\documentclass[11pt,a4paper,oneside]{report}

%---------------------------------------------------------------------
% Encodage & langue
%---------------------------------------------------------------------
\usepackage[utf8]{inputenc}
\usepackage[T1]{fontenc}
\usepackage{lmodern}
\usepackage[french]{babel}
\usepackage{microtype}

%---------------------------------------------------------------------
% Mise en page
%---------------------------------------------------------------------
\usepackage[a4paper,margin=2.5cm,headheight=15pt]{geometry}
\usepackage{graphicx}
\usepackage{xcolor}
\usepackage{booktabs}
\usepackage{longtable}
\usepackage{array}
\usepackage{enumitem}
\usepackage{caption}
\usepackage{fancyhdr}
\usepackage{titlesec}
\usepackage[most]{tcolorbox}
\usepackage{listings}
\usepackage{hyperref}

%---------------------------------------------------------------------
% Charte de couleurs (bleu marine du cahier)
%---------------------------------------------------------------------
\definecolor{rtnavy}{HTML}{1F3A5F}
\definecolor{rtgray}{HTML}{6B7280}
\definecolor{rtlight}{HTML}{EEF2F7}
\definecolor{rtcode}{HTML}{0B1021}
\definecolor{rtaccent}{HTML}{2563EB}

\hypersetup{
  colorlinks=true,
  linkcolor=rtnavy,
  urlcolor=rtaccent,
  citecolor=rtnavy,
  pdftitle={RT-Comops Backend — Rapport d'Analyse et Conception — auth-core},
  pdfauthor={Tchoutzine}
}

%---------------------------------------------------------------------
% En-têtes / pieds de page
%---------------------------------------------------------------------
\pagestyle{fancy}
\fancyhf{}
\fancyhead[L]{\small\itshape\leftmark}
\fancyhead[R]{\small\thepage}
\fancyfoot[C]{\small\color{rtgray}RT-Comops — Rapport d'Analyse \& Conception — \texttt{auth-core}}
\renewcommand{\headrulewidth}{0.4pt}
\renewcommand{\footrulewidth}{0.4pt}

%---------------------------------------------------------------------
% Titres de chapitres / sections
%---------------------------------------------------------------------
\titleformat{\chapter}[display]
  {\normalfont\huge\bfseries\color{rtnavy}}
  {\filright\Large Chapitre~\thechapter}{0.6ex}{\Huge}
\titleformat{\section}
  {\normalfont\Large\bfseries\color{rtnavy}}{\thesection}{1em}{}
\titleformat{\subsection}
  {\normalfont\large\bfseries\color{rtnavy!85}}{\thesubsection}{1em}{}

%---------------------------------------------------------------------
% Boîtes thématiques
%---------------------------------------------------------------------
\newtcolorbox{rolebox}[1][]{
  colback=rtlight, colframe=rtnavy, boxrule=0.8pt, arc=2pt,
  left=8pt,right=8pt,top=6pt,bottom=6pt, #1}

\newtcolorbox{invbox}[1][Invariants]{
  enhanced, colback=white, colframe=rtnavy!70, boxrule=0.6pt, arc=1pt,
  title=#1, coltitle=white, colbacktitle=rtnavy,
  fonttitle=\bfseries, left=8pt,right=8pt,top=6pt,bottom=6pt}

\newtcolorbox{notebox}{
  colback=rtlight!60, colframe=rtaccent, boxrule=0pt,
  leftrule=3pt, arc=0pt, left=8pt,right=8pt,top=5pt,bottom=5pt}

%---------------------------------------------------------------------
% Listings (JSON / bash)
%---------------------------------------------------------------------
\lstdefinestyle{shell}{
  backgroundcolor=\color{rtcode}, basicstyle=\ttfamily\scriptsize\color{white},
  breaklines=true, columns=fullflexible, frame=none,
  keywordstyle=\color{cyan}, stringstyle=\color{green!60!white},
  showstringspaces=false, xleftmargin=6pt, xrightmargin=6pt,
  framexleftmargin=6pt, aboveskip=8pt, belowskip=8pt}
\lstset{style=shell}

\newcommand{\code}[1]{\texttt{\small #1}}
\newcommand{\httpget}{\textcolor{white}{\colorbox{rtaccent}{\scriptsize\ttfamily\,GET\,}}}
\newcommand{\httppost}{\textcolor{white}{\colorbox{green!45!black}{\scriptsize\ttfamily\,POST\,}}}
\newcommand{\httpput}{\textcolor{white}{\colorbox{orange!80!black}{\scriptsize\ttfamily\,PUT\,}}}

%=====================================================================
\begin{document}
%=====================================================================

%--------------------------- Page de titre ---------------------------
\begin{titlepage}
\centering
\vspace*{2.2cm}
{\color{rtnavy}\rule{\linewidth}{2pt}}\\[0.6cm]
{\Huge\bfseries\color{rtnavy} RT-Comops Backend\par}
\vspace{0.4cm}
{\LARGE Rapport d'Analyse \& Conception\par}
\vspace{0.2cm}
{\Large\itshape Module \texttt{auth-core} — YowAuth\par}
\vspace{0.3cm}
{\large Fournisseur d'identité (IdP) central — ConceptionKernel V3\par}
\vspace{0.8cm}
{\color{rtnavy}\rule{\linewidth}{0.8pt}}\\[1.2cm]

\begin{tabular}{rl}
\textbf{Module}      & Identité \& Accès — \texttt{auth-core}\\[2pt]
\textbf{Auteur}      & Tchoutzine\\[2pt]
\textbf{Cours}       & Réseaux Intelligents — Master\\[2pt]
\textbf{Périmètre}   & Comptes utilisateurs, authentification, OIDC\\[2pt]
\textbf{Signature}   & JWT \textbf{RS256} — JWKS public\\[2pt]
\end{tabular}

\vfill
{\large\bfseries Cahier d'Analyse Technique\par}
\vspace{0.4cm}
{\color{rtgray}\small Document de référence pour la conception, le développement et l'évolution du module d'authentification.}
\vspace{1.5cm}
\end{titlepage}

%--------------------------- Table des matières ----------------------
\pagenumbering{roman}
\tableofcontents
\cleardoublepage
\pagenumbering{arabic}

%=====================================================================
\chapter{Introduction}
%=====================================================================
\section{Objet du document}
Ce rapport décrit l'analyse et la conception du module \textbf{\code{auth-core}} (nommé
\emph{YowAuth}) de la plateforme RT-Comops Backend. Il constitue le document de référence
technique et fonctionnelle du sous-système d'authentification : modèle de domaine, ports et
adaptateurs, cas d'usage, diagrammes de séquence, invariants métier, surface d'API exposée,
mécanismes de sécurité (JWT RS256, JWKS, OIDC) et stratégie de vérification.

Le document suit les mêmes conventions que le cahier d'analyse global RT-Comops : pour le
core étudié sont fournis la description détaillée, les rôles primaire et secondaires, le
diagramme de classes, le diagramme des cas d'utilisation et un ensemble de diagrammes de
séquence couvrant les interactions techniques les plus significatives.

\section{Positionnement d'auth-core dans le kernel}
La plateforme est un \textbf{monolithe modulaire à architecture hexagonale}, organisé en
une vingtaine de \emph{cores} fonctionnels plus un module \code{bootstrap} exécutable.
\code{auth-core} appartient au domaine \emph{Identité \& Accès}, aux côtés de
\code{actor-core} (identités humaines) et \code{roles-core} (RBAC).

\begin{rolebox}
\code{auth-core} sépare clairement l'\textbf{identité humaine} (\code{Actor}, gérée par
\code{actor-core}) du \textbf{compte technique} (\code{UserAccount}) utilisé pour
s'authentifier et porter un contexte utilisateur dans les appels au kernel. La signature
RS256 des JWT et l'exposition du JWKS restent portées par \code{kernel-core} ;
\code{auth-core} orchestre l'émission, la sélection de contexte, les challenges de sécurité
et l'exposition OIDC après authentification réussie et résolution des permissions.
\end{rolebox}

\section{Périmètre fonctionnel}
Le module couvre le cycle IAM complet attendu d'un fournisseur d'identité first-party :
\begin{itemize}[leftmargin=1.4em,itemsep=1pt]
  \item création et cycle de vie des comptes \code{UserAccount} ;
  \item authentification locale tenant-scoped et découverte globale de contextes ;
  \item sélection de contexte tenant/organisation ;
  \item émission d'un \textbf{JWT d'accès RS256} enrichi des permissions ;
  \item self-service utilisateur (\code{/users/me}, plan, onboarding) ;
  \item flux de challenge et durcissement (CAPTCHA, OTP, MFA, vérification email/téléphone) ;
  \item réinitialisation de mot de passe (SMTP ou mode \code{PREVIEW\_ONLY}) ;
  \item provider \textbf{OIDC/OAuth2} central (métadonnées, token-exchange, userinfo, introspection, JWKS) ;
  \item projection des \textbf{organisations accessibles} et de leurs services effectifs.
\end{itemize}

%=====================================================================
\chapter{Description et rôles du core}
%=====================================================================
\section{Description du core}
L'\code{auth-core} est le module responsable des comptes utilisateurs et de
l'authentification locale. Il expose désormais un flux HTTP en variantes complémentaires
couvrant le cycle complet IAM :
\begin{itemize}[leftmargin=1.4em,itemsep=1pt]
  \item le flux historique \httppost~\code{/api/auth/login}, tenant-scoped, qui suppose un \code{X-Tenant-Id} connu ;
  \item le flux \emph{email-first} \httppost~\code{/api/auth/identify}, qui détermine si le principal doit poursuivre vers \code{SIGN\_IN\_PASSWORD} ou \code{SIGN\_UP} ;
  \item le flux moderne \httppost~\code{/api/auth/discover-contexts} puis \httppost~\code{/api/auth/select-context}, permettant un login global par principal/password avant sélection du contexte tenant/organisation ;
  \item les flux de self-service : \code{discover-sign-up-contexts}, \code{sign-up}, \code{forgot-password}, \code{password-reset/issue}, \code{reset-password}, \code{email-verification/*} ;
  \item les flux de challenge : \code{captcha}, \code{otp}, \code{phone-verification/*}, \code{mfa/*}, \code{login/mfa/confirm} ;
  \item le provider OIDC central : \code{/.well-known/openid-configuration}, \code{/.well-known/oauth-authorization-server}, \code{/oauth2/token}, \code{/oauth2/userinfo}, \code{/oauth2/introspect}, \code{/.well-known/jwks.json}.
\end{itemize}

\begin{notebox}
\textbf{Note.} Le backend ne possède pas encore de catalogue métier des tenants avec nom
d'affichage. Le flux de découverte de contexte retourne donc des \code{tenantId} techniques,
accompagnés des organisations accessibles dans chacun de ces contextes.
\end{notebox}

\section{Rôle principal}
\begin{rolebox}
\textbf{Rôle principal :} gérer le cycle de vie des comptes \code{UserAccount}, assurer
l'authentification locale dans un contexte tenant donné ou via une découverte globale de
contextes, puis retourner un JWT d'accès signé contenant l'identité utilisateur, le contexte
sélectionné et les permissions résolues.
\end{rolebox}

\section{Rôles secondaires}
\begin{enumerate}[leftmargin=1.6em,itemsep=2pt]
  \item \textbf{Création contrôlée des comptes.} \httppost~\code{/api/auth/register} est
  réservé aux appels authentifiés portés par une \code{ClientApplication} et un utilisateur
  disposant des permissions IAM (\code{system:admin} / \code{iam:admin}). Les identifiants
  sont normalisés et le mot de passe haché (\code{BCryptPasswordEncoder}) avant persistance.
  \item \textbf{Authentification locale.} Deux parcours : \code{/login} (résolution d'un
  compte par principal dans un tenant, puis vérification du mot de passe) et
  \code{/discover-contexts} (résolution globale, filtrage par mot de passe valide, émission
  d'un jeton court de sélection de contexte).
  \item \textbf{Sélection de contexte tenant/organisation.} \code{/select-context} consomme
  le jeton court, valide le \code{contextId} et l'appartenance à l'organisation, puis émet le
  JWT final portant \code{tenantId} et éventuellement \code{organizationId}.
  \item \textbf{Émission d'un JWT d'accès RS256.} Après authentification, \code{auth-core}
  résout les permissions puis délègue au \code{UserSessionTokenService} de \code{kernel-core}
  l'émission du bearer token RS256.
  \item \textbf{Auto-service utilisateur.} \code{GET /api/users/me}, \code{PUT
  /users/me/plan}, \code{PUT /users/me/onboarding}, \code{PUT /users/me/identity-onboarding}.
  Le profil retourne aussi \code{organizations[]} et, pour chaque organisation, les services
  effectivement disponibles.
  \item \textbf{Projection des entitlements organisationnels.} \code{auth-core} interroge
  \code{organization-core} pour agréger les organisations accessibles et leurs services
  effectifs, permettant au frontend d'afficher/masquer les modules métier.
  \item \textbf{Compatibilité ascendante.} Le flux historique \code{/login} est conservé pour
  les consommateurs connaissant déjà le \code{tenantId}.
  \item \textbf{Journalisation applicative.} Les opérations sensibles
  (\code{USER\_REGISTERED}, \code{USER\_SIGNED\_UP}, \code{USER\_LOGIN},
  \code{USER\_PASSWORD\_RESET\_*}, \code{USER\_EMAIL\_VERIFICATION\_*},
  \code{USER\_PLAN\_UPDATED}, \code{USER\_ONBOARDING\_UPDATED}) sont inscrites dans l'audit
  système via \code{kernel-core}.
  \item \textbf{Inscription contextuelle sans tenant explicite.} Via
  \code{discover-sign-up-contexts} avec un \code{organizationCode} : le résultat inclut un
  \code{selectionToken} signé et un ou plusieurs \code{contextId} réutilisables par
  \code{sign-up}.
  \item \textbf{Réinitialisation de mot de passe.} \code{forgot-password} résout tous les
  comptes correspondant à un principal, émet un jeton court de sélection ; puis
  \code{password-reset/issue} produit un jeton signé de reset. Si un provider SMTP est
  configuré, l'email est envoyé ; sinon, mode explicite \code{PREVIEW\_ONLY}.
  \item \textbf{Vérification d'adresse email.} \code{email-verification/request} émet un jeton
  signé associé à l'utilisateur ; \code{email-verification/confirm} valide ce jeton et
  renseigne \code{emailVerifiedAt}.
  \item \textbf{OTP, téléphone et MFA.} Challenges OTP (email / SMS / WhatsApp),
  \code{phone-verification/*}, activation MFA en deux temps puis exigence de
  \code{login/mfa/confirm} aux connexions suivantes.
  \item \textbf{CAPTCHA et anti-abus.} \code{captcha} génère un challenge court,
  \code{captcha/verify} produit un jeton réutilisable par \code{sign-up}.
  \item \textbf{Provider OIDC central.} Métadonnées OIDC/OAuth2, échange de jeton
  (\code{token-exchange}), \code{userinfo}, \code{introspect}, adossés au JWKS RS256 du kernel.
\end{enumerate}

%=====================================================================
\chapter{Architecture hexagonale du module}
%=====================================================================
\section{Structure interne}
Le module suit la structure hexagonale canonique. Le domaine ne connaît ni le web ni la
base : il ne parle qu'à des interfaces (ports).

\begin{center}
\begin{tabular}{@{}p{5.2cm}p{8.6cm}@{}}
\toprule
\textbf{Couche / package} & \textbf{Contenu}\\
\midrule
\code{domain.model} & Agrégat \code{UserAccount}, invariants métier, énumérations de statut.\\
\code{application.port.in} & Cas d'usage entrants (\code{LoginUseCase}, \code{RegisterUserUseCase}, \code{PublicSignUpUseCase}, \code{DiscoverLoginContextsUseCase}, \dots).\\
\code{application.port.out} & Ports sortants : \code{UserAccountRepository}, \code{RefreshTokenRepository}, \code{UserOrganizationAccessDirectory}, \code{SignUpContextDirectory}.\\
\code{application.service} & Orchestration : \code{AuthApplicationService}, \code{AuthContextApplicationService}, \code{AuthOidcService}, \code{AuthSharedSessionService} et les services de jetons.\\
\code{adapter.in.web} & Contrôleurs WebFlux (\code{AuthController}, \code{AuthOidcController}, \code{SessionTokensController}) et DTO.\\
\code{adapter.out.persistence} & \code{UserAccountR2dbcRepositoryAdapter}, \code{InMemoryUserAccountRepository} (\code{@Profile("test-memory")}).\\
\code{config} & Assemblage Spring, \code{AuthModuleConfiguration}, propriétés (\code{AuthSsoProperties}, \code{AuthEmailProperties}, \code{LoginThrottleProperties}, \code{RefreshTokenProperties}).\\
\bottomrule
\end{tabular}
\end{center}

\section{Dépendances et ports sortants}
\code{auth-core} dépend au compile de \code{common-core}, \code{kernel-core} et
\code{actor-core}. Trois ports sortants matérialisent les couplages inter-cores, résolus par
des adaptateurs alternatifs selon l'assemblage :

\begin{center}
\begin{tabular}{@{}p{5.4cm}p{8.4cm}@{}}
\toprule
\textbf{Port sortant} & \textbf{Rôle / implémentation}\\
\midrule
\code{UserAccountRepository} & Persistance des comptes (R2DBC en prod, in-memory en test).\\
\code{RefreshTokenRepository} & Persistance des refresh tokens.\\
\code{UserOrganizationAccess\allowbreak Directory} & Liste des organisations accessibles d'un user ; adaptateur réel vers \code{organization-core} (via \code{bootstrap}).\\
\code{SignUpContextDirectory} & Contextes d'inscription par code organisation ; adaptateur réel vers \code{organization-core}.\\
\bottomrule
\end{tabular}
\end{center}

\begin{notebox}
\textbf{Extensibilité.} Les trois ports liés à l'organisation peuvent être servis par des
adaptateurs de sortie \emph{alternatifs} (par exemple in-memory seedés) sans modifier le
domaine — propriété exploitée pour extraire et déployer \code{auth-core} de façon isolée
(chapitre~\ref{chap:poc}).
\end{notebox}

%=====================================================================
\chapter{Modèle de domaine}
%=====================================================================
\section{Diagramme de classes}
\begin{figure}[h]
  \centering
  % Rendre diagrams/auth-core-class.puml en PNG (voir README) et décommenter :
  % \includegraphics[width=\linewidth]{figures/auth-core-class.png}
  \fbox{\parbox{0.95\linewidth}{\centering\vspace{1cm}\itshape\color{rtgray}
  Diagramme de classes — \code{auth-core}\\[4pt]
  (image générée depuis \texttt{diagrams/auth-core-class.puml})\vspace{1cm}}}
  \caption{Diagramme de classes du module auth-core.}
\end{figure}

\subsection{Agrégat \code{UserAccount}}
Modèle de domaine immuable représentant un compte utilisateur. Attributs principaux :
\code{id}, \code{tenantId}, \code{actorId}, \code{username}, \code{email},
\code{phoneNumber}, \code{passwordHash}, \code{authProvider}, \code{externalSubject},
\code{status}, \code{plan}, \code{onboardingStatus}, \code{onboardingStep},
\code{accountType}, \code{businessType}, \code{onboardingPayload}, \code{emailVerified},
\code{emailVerifiedAt}, \code{phoneVerified}, \code{phoneVerifiedAt}, \code{mfaEnabled},
\code{mfaChannel}. Chaque mutation retourne un nouvel objet (immuabilité, sûre en réactif) :
\code{register()}, \code{markEmailVerified()}, \code{enableMfa()}, \code{disableMfa()},
\code{updatePlan()}, \code{updateOnboarding()}, \code{linkExternalIdentity()},
\code{updatePassword()}, \code{markPhoneVerified()}.

\subsection{Services de jetons}
\begin{center}
\begin{tabular}{@{}p{6.2cm}p{7.6cm}@{}}
\toprule
\textbf{Service} & \textbf{Rôle}\\
\midrule
\code{AuthSsoSessionTokenService} & « Passeport » SSO après login (contextes multi-tenant embarqués).\\
\code{AuthSharedSessionService} & Session partagée multi-tenant ; \code{issueForUser()}, \code{userInfo()}.\\
\code{AuthContextSelectionToken\allowbreak Service} & Jeton court et temporaire de sélection de contexte.\\
\code{AuthPasswordResetToken\allowbreak Service} & Jetons de réinitialisation de mot de passe.\\
\code{AuthEmailVerificationToken\allowbreak Service} & Jetons de vérification d'email (TTL 24h).\\
\code{AuthSignUpSelectionToken\allowbreak Service} & Jetons de sélection de contexte d'inscription.\\
\code{AuthChallengeTokenService} & OTP et CAPTCHA.\\
\code{JwtTokenService} (kernel) & Signe/vérifie les JWT RS256 ; charge une clé PEM ou auto-génère une paire RSA en mémoire ; exporte le JWKS public.\\
\bottomrule
\end{tabular}
\end{center}

\section{Invariants métier}
\begin{invbox}[Invariants du auth-core]
\begin{itemize}[leftmargin=1.2em,itemsep=1pt]
  \item \code{username} est unique par tenant et normalisé en minuscule.
  \item \code{email} est unique par tenant et normalisé en minuscule.
  \item \code{passwordHash} est obligatoire pour le flux local courant.
  \item \code{authProvider} est obligatoire et décrit l'origine du compte.
  \item \code{plan} et \code{onboardingStatus} ne peuvent pas être vides ; \code{onboardingStep} $\geq 0$.
  \item \code{emailVerifiedAt} / \code{phoneVerifiedAt} nuls signifient non confirmés.
  \item Un login MFA-enabled ne doit pas retourner de session finale avant validation du challenge MFA.
  \item OTP, CAPTCHA, jetons de reset, de sélection et MFA sont courts, signés et bornés dans le temps.
  \item Le jeton court de sélection de contexte est signé, de courte durée, et n'est pas un bearer token métier.
  \item Le jeton de reset et le jeton de vérification email sont signés par \code{kernel-core} via la même infrastructure JWT RS256.
  \item Le provider OIDC central ne remplace pas le RBAC métier ; il publie des tokens alignés avec le contexte sélectionné.
  \item Les réponses \code{login}, \code{select-context} et \code{/users/me} doivent refléter les organisations réellement accessibles et leurs services effectifs.
\end{itemize}
\end{invbox}

%=====================================================================
\chapter{Cas d'utilisation}
%=====================================================================
\begin{figure}[h]
  \centering
  % \includegraphics[width=0.72\linewidth]{figures/auth-core-usecase.png}
  \fbox{\parbox{0.9\linewidth}{\centering\vspace{0.8cm}\itshape\color{rtgray}
  Diagramme des cas d'utilisation — \code{auth-core}\\[4pt]
  (image générée depuis \texttt{diagrams/auth-core-usecase.puml})\vspace{0.8cm}}}
  \caption{Cas d'utilisation du module auth-core.}
\end{figure}

\noindent\textbf{Acteurs.}
\begin{itemize}[leftmargin=1.4em,itemsep=1pt]
  \item \emph{Administrateur IAM} — crée un compte utilisateur (\code{register}).
  \item \emph{Utilisateur} — consulte/ajuste son profil, son plan, son onboarding ; émet des challenges de vérification.
  \item \emph{Backend consommateur} — identifie un principal, authentifie, découvre et sélectionne un contexte, inscrit en self-service, réinitialise le mot de passe.
\end{itemize}
Le cas transverse \emph{« Journaliser les opérations de sécurité »} est inclus (\code{include})
par la plupart des cas d'usage ; le cas \emph{« Émettre un JWT d'accès RS256 »} est inclus par
les flux d'authentification aboutis.

%=====================================================================
\chapter{Surface d'API exposée}
%=====================================================================
Toutes les routes \code{/api/**} sont soumises à la chaîne de sécurité du kernel :
\code{ClientApplication} valide (\code{X-Client-Id} + \code{X-Api-Key}), autorisée sur le
service, quota backend, abonnement d'organisation puis permission utilisateur au bon scope.
Les réponses suivent l'enveloppe \code{ApiResponse<T> = \{success, data, message, errorCode, timestamp\}}.

\section{Authentification et contextes}
\begin{center}\small
\begin{longtable}{@{}p{1.2cm}p{5.9cm}p{6.2cm}@{}}
\toprule
\textbf{M.} & \textbf{Chemin} & \textbf{Corps / rôle}\\
\midrule\endhead
POST & \code{/api/auth/sign-up} & \code{tenantId, firstName, lastName, username, email, password, \dots}\\
POST & \code{/api/auth/login} & \code{\{principal, password\}} → accessToken, sessionToken, sharedSession, organisations\\
POST & \code{/api/auth/identify} & \code{\{principal\}} — SIGN\_UP vs SIGN\_IN\\
POST & \code{/api/auth/register} & (admin) \code{username, email, authProvider, \dots}\\
POST & \code{/api/auth/discover-contexts} & \code{\{principal, password\}} → selectionToken + contexts\\
POST & \code{/api/auth/select-context} & \code{\{selectionToken, contextId, organizationId?\}}\\
POST & \code{/api/auth/discover-sign-up-contexts} & \code{\{organizationCode\}}\\
POST & \code{/api/auth/refresh} & \code{\{refreshToken\}}\\
POST & \code{/api/auth/logout} & invalide la session courante\\
\bottomrule
\end{longtable}
\end{center}

\section{Vérification, mot de passe, challenges}
\begin{center}\small
\begin{longtable}{@{}p{1.2cm}p{6.6cm}p{5.5cm}@{}}
\toprule
\textbf{M.} & \textbf{Chemin} & \textbf{Corps / rôle}\\
\midrule\endhead
POST & \code{/api/auth/email-verification/request} & jeton pour l'utilisateur courant\\
POST & \code{/api/auth/email-verification/resend} & \code{\{principal\}} (mode PREVIEW : token en clair)\\
POST & \code{/api/auth/email-verification/confirm} & \code{\{verificationToken\}}\\
POST & \code{/api/auth/forgot-password} & \code{\{principal\}} → selectionToken\\
POST & \code{/api/auth/password-reset/issue} & \code{\{selectionToken, contextId\}}\\
POST & \code{/api/auth/reset-password} & \code{\{resetToken, newPassword\}}\\
POST & \code{/api/auth/otp} , \code{/otp/verify} & émission / vérification OTP\\
POST & \code{/api/auth/captcha} , \code{/captcha/verify} & challenge / jeton CAPTCHA\\
POST & \code{/api/auth/mfa/enable|confirm|disable} & activation MFA en deux temps\\
POST & \code{/api/auth/login/mfa/confirm} & \code{\{mfaToken, code\}}\\
POST & \code{/api/auth/phone-verification/request|confirm} & vérification téléphone\\
\bottomrule
\end{longtable}
\end{center}

\section{Self-service et OIDC}
\begin{center}\small
\begin{longtable}{@{}p{1.2cm}p{6.6cm}p{5.5cm}@{}}
\toprule
\textbf{M.} & \textbf{Chemin} & \textbf{Rôle}\\
\midrule\endhead
GET & \code{/api/users/me} & profil + organisations accessibles\\
PUT & \code{/api/users/me/plan} & met à jour le plan\\
PUT & \code{/api/users/me/onboarding} & met à jour l'onboarding\\
PUT & \code{/api/users/me/identity-onboarding} & type de compte / activité / payload\\
GET & \code{/.well-known/openid-configuration} & découverte OIDC\\
GET & \code{/.well-known/oauth-authorization-server} & métadonnées OAuth2\\
GET & \code{/.well-known/jwks.json} & clé publique RSA (RS256)\\
POST & \code{/oauth2/token} & token-exchange (\code{context\_id} + \code{service\_code})\\
GET/POST & \code{/oauth2/userinfo} & infos utilisateur depuis un access token\\
POST & \code{/oauth2/introspect} & introspection de jeton\\
\bottomrule
\end{longtable}
\end{center}

%=====================================================================
\chapter{Sécurité et cryptographie}
%=====================================================================
\section{Signature asymétrique RS256}
Les JWT d'accès sont signés en \textbf{RS256} : \code{auth-core} (via le
\code{JwtTokenService} du kernel) signe avec la \textbf{clé privée} RSA ; tout service
vérifie avec la \textbf{clé publique}. Aucun secret partagé n'est distribué : chaque
microservice vérifie un jeton hors-ligne. La clé publique est exposée au format \textbf{JWKS}
sur \code{/.well-known/jwks.json}.

\begin{lstlisting}[caption={Exemple de payload d'access token (claims décodés).}]
{
  "sub": "c256dcb1-...",   "aud": "iwm-api",   "iss": "iwm-backend",
  "svc": "SALES",          "azp": "poc-client", "sso": true,
  "tid": "00000000-0000-0000-0000-000000000001",
  "actor": "bafb4f41-...", "permissions": [], "exp": 1782648517
}
\end{lstlisting}

\section{Mode de livraison PREVIEW\_ONLY}
Lorsqu'aucun provider SMTP n'est configuré (\code{iwm.auth.email.enabled=false} ou absence
de \code{from}), \code{AuthEmailDeliveryService} bascule en mode \code{PREVIEW\_ONLY} et
renvoie le jeton de challenge \textbf{en clair} dans la réponse (\code{challengeTokenPreview}).
Ce mode permet d'exercer les flux de vérification email et de reset sans serveur mail — utile
en dev, en test et pour un POC. En production, la livraison est réellement SMTP
(\code{deliveryMode = SMTP}).

\section{Mode strict de connexion}
Un compte \code{LOCAL} dont l'email n'est pas vérifié ne peut pas se connecter : le
\code{sign-up} renvoie \code{EMAIL\_VERIFICATION\_REQUIRED} sans session tant que l'email
n'est pas confirmé.

%=====================================================================
\chapter{Diagrammes de séquence (DS-AU)}
%=====================================================================
Cette section reprend les scénarios d'interaction les plus significatifs. Chaque diagramme
possède une source PlantUML dans \texttt{diagrams/}.

\newcommand{\dsfig}[2]{%
  \begin{figure}[h]\centering
  % \includegraphics[width=\linewidth]{figures/#1.png}
  \fbox{\parbox{0.95\linewidth}{\centering\vspace{0.6cm}\itshape\color{rtgray}
  #2\\[3pt](\texttt{diagrams/#1.puml})\vspace{0.6cm}}}
  \caption{#2}\end{figure}}

\section{DS-AU-01 — Création d'un compte par un administrateur IAM}
Requête \httppost~\code{/api/auth/register} portée par une \code{ClientApplication} et un
bearer admin. Le service vérifie l'unicité \code{username}, encode le mot de passe, persiste
le \code{UserAccount} et journalise \code{USER\_REGISTERED}.
\dsfig{ds-au-01-register}{DS-AU-01 — Création d'un compte utilisateur (admin IAM)}

\section{DS-AU-02 / 03 — Authentification tenant-scoped et émission du JWT}
\code{/login} résout le compte par principal dans le tenant, vérifie le mot de passe
(BCrypt), résout les permissions (\code{ReactivePermissionResolver}), puis émet l'access
token RS256 via \code{UserSessionTokenService} et le retourne avec les organisations.
\dsfig{ds-au-02-login}{DS-AU-02/03 — Login tenant-scoped puis émission du JWT RS256}

\section{DS-AU-07 / 08 — Découverte et sélection de contexte}
\code{/discover-contexts} résout globalement tous les comptes correspondant au principal,
filtre par mot de passe valide, agrège les organisations (\code{UserOrganizationAccessDirectory})
et émet un jeton court de sélection. \code{/select-context} le consomme, valide le contexte et
émet le JWT final contextualisé.
\dsfig{ds-au-07-context}{DS-AU-07/08 — Découverte globale puis sélection de contexte}

\section{DS-AU-10 — Réinitialisation de mot de passe (SMTP ou PREVIEW\_ONLY)}
\code{forgot-password} → jeton de sélection ; \code{password-reset/issue} → jeton de reset
signé ; livraison SMTP ou repli \code{PREVIEW\_ONLY} ; \code{reset-password} vérifie le jeton
et met à jour le \code{passwordHash}.
\dsfig{ds-au-10-reset}{DS-AU-10 — Reset de mot de passe (livraison SMTP ou PREVIEW\_ONLY)}

\section{DS-AU-11 — Vérification d'email après authentification}
\code{email-verification/request} émet un jeton signé ; \code{confirm} le valide et
renseigne \code{emailVerifiedAt}.
\dsfig{ds-au-11-emailverif}{DS-AU-11 — Vérification d'email}

%=====================================================================
\chapter{Flow OIDC / token-exchange}
%=====================================================================
Le provider OIDC central expose la découverte et le grant \code{token-exchange} (RFC 8693) :
un jeton de session SSO est échangé, avec le \code{context\_id} (organisation) et le
\code{service\_code}, contre un access token dédié au service. Le \code{userinfo} restitue
l'identité depuis ce token.

\begin{lstlisting}[caption={Échange token-exchange puis userinfo.}]
curl -X POST $BASE/oauth2/token \
  -H "Content-Type: application/x-www-form-urlencoded" \
  --data-urlencode "grant_type=urn:ietf:params:oauth:grant-type:token-exchange" \
  --data-urlencode "subject_token=$SSO_TOKEN" \
  --data-urlencode "subject_token_type=urn:ietf:params:oauth:token-type:jwt" \
  --data-urlencode "context_id=$CTX" --data-urlencode "service_code=SALES" \
  --data-urlencode "client_id=$CLIENT_ID" --data-urlencode "client_secret=$SECRET"
# => { "access_token": "...RS256...", "issued_token_type": "...access_token" }

curl $BASE/oauth2/userinfo -H "Authorization: Bearer $ACCESS_TOKEN"
# => { "sub", "email", "tenantId", "actorId", "clientId", "serviceCode", "sso" }
\end{lstlisting}

\begin{notebox}
\textbf{Périmètre OIDC.} Le seul \code{grant\_types\_supported} est le \code{token-exchange}.
Il n'existe pas d'\code{authorization\_endpoint} ni d'écran de consentement « à la Google » :
\code{auth-core} est l'IdP des applications \emph{first-party}, pas un fournisseur OAuth pour
applications tierces.
\end{notebox}

%=====================================================================
\chapter{Stratégie de test et couverture}
%=====================================================================
\section{Tests unitaires implémentés}
Le module a été couvert par une suite de \textbf{tests unitaires} exécutés sans base de
données ni contexte Spring (en mémoire) :
\begin{center}\small
\begin{tabular}{@{}p{7.4cm}p{2.0cm}p{4.2cm}@{}}
\toprule
\textbf{Fichier de test} & \textbf{\#} & \textbf{Cible}\\
\midrule
\code{UserAccountTest} & 21 & modèle de domaine, invariants, immuabilité\\
\code{AuthApplicationServiceTest} & 15 & register / login / identify, isolation tenant\\
\code{AuthContextApplicationServiceTest} & 10 & discover / select multi-tenant\\
\code{JwtTokenServiceTest} & 11 & signature RS256, claims, JWKS\\
\code{AuthSsoSessionTokenServiceTest} & 10 & passeport SSO, TTL, rejets\\
\code{AuthSharedSessionServiceTest} & 6 & issueForUser / userInfo\\
\code{AuthOidcServiceTest} & 11 & token-exchange, openid-configuration\\
\midrule
\textbf{Total auth-core} & \textbf{84} & tous verts (BUILD SUCCESS)\\
\code{roles-core} (bonus) & 20 & Role, RBAC scopeé\\
\bottomrule
\end{tabular}
\end{center}

\section{Points de vérification clés}
\begin{itemize}[leftmargin=1.4em,itemsep=1pt]
  \item hachage BCrypt, rejet des doublons \code{username}/\code{email}, isolation par tenant ;
  \item émission d'un token SSO avec contextes embarqués et respect du TTL ;
  \item JWKS exporté contenant une clé publique RSA (sans la clé privée) ;
  \item rejet des tokens corrompus, nuls, de mauvais type ou d'un autre issuer ;
  \item token-exchange : succès \code{Bearer}, claims corrects, rejets \code{invalid\_grant} / \code{invalid\_client} / \code{invalid\_request}.
\end{itemize}

\section{Couverture restant à compléter}
Flux secondaires documentés mais non encore couverts par des tests dédiés côté serveur : mot de passe
oublié (bout en bout), inscription avec sélection de contexte,
\code{AuthChallengeTokenService} (CAPTCHA), tests d'intégration R2DBC/HTTP globaux, et cas
limites (compte \code{SUSPENDED}/\code{REVOKED}, token expiré). Les flux de double authentification (MFA),
d'activation/désactivation d'OTP et de modification de forfait ont été quant à eux entièrement
implémentés et validés fonctionnellement via le client d'interface utilisateur POC (Next.js).

%=====================================================================
\chapter{Déploiement isolé et démonstration (POC)}\label{chap:poc}
%=====================================================================
Grâce à l'architecture hexagonale, \code{auth-core} a été extrait dans un \emph{slice}
autonome et déployé indépendamment du reste du monolithe. Le slice embarque
\code{common-core}, \code{kernel-core}, \code{file-core}, \code{actor-core},
\code{roles-core} et \code{auth-core}, plus un runner Spring Boot exécutable. Les trois ports
liés à l'organisation sont servis par des adaptateurs \textbf{in-memory seedés}, et la
persistance est en mémoire (profil \code{test-memory}, JWT RSA auto-généré).

\section{Interface Client de Démonstration (Next.js)}
Afin de valider et présenter le rôle de vitrine de YowAuth, une application web
\textbf{Next.js} moderne et responsive a été développée et intégrée avec le Kernel.
Elle implémente les flux IHM suivants dans le respect des règles ergonomiques :

\subsection{Connexion en 2 étapes et gestion du MFA}
L'interface de connexion sépare la phase de détection de compte de la validation des secrets :
\begin{enumerate}
  \item \textbf{Étape d'identification} : appel à \code{/api/auth/identify} pour vérifier l'existence de l'identifiant.
  \item \textbf{Saisie du mot de passe} : si le compte existe, saisie du secret. En cas de défi MFA retourné par le serveur (code statut \code{202 Accepted}), l'interface bascule dynamiquement sur l'écran OTP.
  \item \textbf{Validation OTP} : saisie du code reçu par e-mail et confirmation finale via \code{/api/auth/login/mfa/confirm}. En mode démo, le code est pré-affiché dans l'interface pour faciliter l'évaluation.
\end{enumerate}

\subsection{Centre de contrôle utilisateur & Sécurité}
Après connexion, l'utilisateur accède à un espace de contrôle structuré en deux colonnes :
\begin{itemize}
  \item \textbf{Gestion de l'abonnement} : mise à jour du plan de l'utilisateur (\code{FREE}, \code{PREMIUM}, \code{ENTERPRISE}) en direct via l'endpoint \code{PUT /api/users/me/plan}.
  \item \textbf{Enrôlement et Révocation MFA} : activation ou désactivation à la volée de la double authentification (\code{enable}, \code{confirm}, \code{disable}) avec validation par code de sécurité.
  \item \textbf{Visualisation des Privilèges} : décodage en temps réel du jeton JWT et affichage des permissions effectives sous forme de badges.
\end{itemize}

\subsection{Simulateur SSO OIDC interactif}
Pour matérialiser le rôle de passerelle de YowAuth, le clic sur une application active (ex: \emph{Ventes}, \emph{Comptabilité}) intercepte la redirection pour lancer une simulation visuelle de Single Sign-On (SSO) :
\begin{enumerate}
  \item Requête de \textbf{Token-Exchange (OIDC RFC 8693)} auprès de \code{/oauth2/token} pour échanger le jeton SSO contre un jeton de service contextualisé à l'organisation active.
  \item Validation du jeton via l'introspection \code{/oauth2/userinfo} pour récupérer le profil utilisateur.
  \item Journalisation colorée de ces étapes techniques dans une console interactive, puis redirection transparente de l'utilisateur vers l'application cible.
\end{enumerate}

\section{Ergonomie et Charte Graphique}
L'interface a été conçue pour offrir une esthétique premium : mise en page en grille responsive à 2 colonnes (applications à gauche, profil et sécurité à droite), gestion adaptative des écrans mobiles, et utilisation d'illustrations représentatives à la place des émojis pour illustrer le catalogue de services de Yowyob.

%=====================================================================
\chapter{Conclusion}
%=====================================================================
\code{auth-core} constitue un fournisseur d'identité complet et cohérent pour l'écosystème
RT-Comops : séparation stricte identité / compte, authentification multi-parcours,
multi-tenant et multi-organisation, émission de JWT \textbf{RS256} vérifiables, self-service
utilisateur, provider OIDC central par \emph{token-exchange}, et une batterie d'invariants de
sécurité. L'architecture hexagonale garantit la testabilité et l'extractibilité du module,
démontrée par un déploiement isolé fonctionnel. Les évolutions naturelles portent sur la
couverture des flux secondaires (reset, MFA, OTP), les tests d'intégration, et — si un usage
tiers l'exige — l'ajout d'un \code{authorization\_endpoint} avec écran de consentement.

\end{document}
